\documentclass[conference]{IEEEtran}
\IEEEoverridecommandlockouts
\usepackage{cite}
\usepackage{amsmath,amssymb,amsfonts}
\usepackage{algorithmic}
\usepackage{graphicx}
\usepackage{textcomp}
\usepackage{xcolor}
\usepackage{booktabs}
\usepackage{multirow}
\usepackage{url}

\def\BibTeX{{\rm B\kern-.05em{\sc i\kern-.025em b}\kern-.08em
    T\kern-.1667em\lower.7ex\hbox{E}\kern-.125emX}}

\begin{document}

\title{Early-Stopping Activation Steering for Factual Preference Alignment in Vietnamese Medical Language Models\\
\thanks{Corresponding Author: Phan Do Thanh Tuan (Email: phandothanhtuan1704@gmail.com).}
}

\author{\IEEEauthorblockN{Phan Do Thanh Tuan}
\IEEEauthorblockA{\textit{Department of Information Technology} \\
\textit{FPT University Gia Lai, Vietnam}\\
Email: phandothanhtuan1704@gmail.com}
}

\maketitle

\begin{abstract}
Small Language Models (SLMs) offer potential for privacy-preserving clinical decision support in local hospital infrastructure. However, when deployed in resource-scarce languages like Vietnamese, SLMs frequently exhibit domain-specific hallucinations---such as miscalculating pediatric dosages, violating pregnancy contraindications, or mistaking drug interaction mechanisms. In this work, we propose an \textbf{Early-Stopping Activation Steering} framework to dynamically shift model outputs toward factual preference at inference time with no model-weight updates. Using our newly constructed \textbf{ViHaluEval-Medical} benchmark (14,700 records derived from the Vietnamese National Drug Formulary, 2nd Edition, 2018), we evaluate steering interventions across Layer~8 of Qwen2.5-7B-Instruct. To prevent activation magnitude inflation caused by continuous steering over long sequences, we introduce a linear decay early-stopping mechanism ($K=16$). Under a dual-evaluation paradigm across a unified question-disjoint test set ($N_{\text{test}}=500$), we report:
\begin{enumerate}
    \item \textbf{High-Cap Natural Completion Evaluation}: Evaluating steering interventions under long-sequence generation (\texttt{max\_new\_tokens=800}), achieving an EOS hit rate range of 99.80\%--100.00\% while demonstrating that Hard Cutoff ($K=16$) achieves the highest reference-preference accuracy of 70.60\% (+5.20~pp over baseline 65.40\%, $p_{\text{adj}}=0.0026$).
    \item \textbf{Controlled Bounded-Generation Evaluation}: Evaluating under \texttt{max\_new\_tokens=200}, early-stopping steering ($\alpha_0=18.0$) increases BERTScore reference-preference accuracy from 73.00\% to 77.20\% (+4.20~pp, $p=0.0055$). Empirical hidden-state activation tracking shows linear decay reduces post-hook norm to 51.20 at $t=50$ (mitigating magnitude growth), while $+v_{\text{steer}}$ outperforms non-oracle BM25 RAG by +8.60~pp with zero context overhead.
\end{enumerate}
\end{abstract}

\begin{IEEEkeywords}
Activation steering, hallucination mitigation, clinical language models, Vietnamese NLP, inference-time intervention, representation engineering.
\end{IEEEkeywords}

\section{Introduction}
Large Language Models (LLMs) have demonstrated capabilities across general healthcare tasks \cite{ref3}. However, in specialized domains requiring high factual precision---such as clinical pharmacology and prescription guidance in non-English languages---standard models suffer from domain-specific hallucinations \cite{ref7}. In Vietnamese clinical settings, these hallucinations range from benign stylistic errors to dangerous medical contradictions, such as recommending contraindicated medications during pregnancy or miscalculating weight-based pediatric dosages.

Traditional mitigation strategies rely on post-hoc supervised fine-tuning (SFT) \cite{ref10} or Retrieval-Augmented Generation (RAG) \cite{ref5}. While RAG provides external knowledge contexts, models may ignore retrieved evidence when internal activation pathways favor non-truthful generation \cite{ref11}, while incurring heavy context window overheads and retrieval latencies. Supervised fine-tuning demands substantial computational resources and risks catastrophic forgetting of general abilities \cite{ref12}.

Recently, Representation Engineering (RepE) \cite{ref1} and Inference-Time Intervention (ITI) \cite{ref2} have emerged as lightweight, non-destructive alternatives. By extracting a truthfulness steering vector $v_{\text{steer}}$ from hidden activation contrasts and injecting it into intermediate layers during generation, activation steering shifts model outputs toward factual accuracy without modifying underlying weights. Related approaches such as Activation Addition (ActAdd) \cite{ref8} and contrastive activation addition \cite{ref9} have demonstrated the viability of linear steering for behavioral control.

However, continuous activation steering ($K=\infty$) can introduce activation magnitude inflation over long output sequences, manifested as repetitive text loops and degraded syntax. In this paper, we present \textbf{Early-Stopping Activation Steering}, a temporal decay mechanism that restricts vector injection to the first $K=16$ generated tokens.

Our main contributions are as follows:
\begin{enumerate}
    \item \textbf{ViHaluEval-Medical Benchmark}: A benchmark of 14,700 expert-structured Vietnamese clinical Q\&A pairs derived from the official Vietnamese National Drug Formulary (2nd Edition, 2018) \cite{ref6}, featuring a question-disjoint test split ($N_{\text{test}}=500$).
    \item \textbf{Dual-Evaluation Paradigm}: Comprehensive evaluation covering both untruncated natural completion (\texttt{max\_new\_tokens=800}, EOS Hit = 99.80\%--100.00\%) and bounded-generation stress testing (\texttt{max\_new\_tokens=200}).
    \item \textbf{Empirical Activation-Norm Dynamics}: Hidden-state activation tracking ($\|\hat{h}_8^{(t)}\|_2$) demonstrating continuous steering norm elevation (56.18 vs 54.21 baseline at $t=50$) and linear decay magnitude relaxation (51.20).
    \item \textbf{Directional Specificity \& Multi-Random Placebo}: Negative steering ($-v_{\text{steer}}$, 62.80\% accuracy) and distributional validation across $N=100$ independent random vectors ($55.82\% \pm 4.15\%$ sample SD) confirming $+v_{\text{steer}}$ ($77.20\%$) operates at a descriptive $+5.15\sigma$ distance above random direction perturbation (empirical randomization $p=1/101=0.0099 < 0.01$), alongside a non-oracle BM25 RAG baseline (68.60\% accuracy, Recall@1 = 19.20\%).
    \item \textbf{Reproducible Statistical Validation}: Exact paired $2\times 2$ binomial McNemar testing ($p=0.0055$, 95\% CI $[+1.37\text{ pp}, +7.03\text{ pp}]$) and Wilcoxon signed-rank testing ($p=0.017$) confirming reference-preference gains.
    \item \textbf{Controlled Factorial Ablation}: Controlled $3\times 3$ equal-initial-$\alpha_0$ matrix alongside a separately specified matched cumulative dose control ($\alpha_{\text{match}}$).
\end{enumerate}

\section{Related Work}

\subsection{Hallucination in Medical Language Models}
LLM hallucinations in clinical domains pose unique risks compared to general-purpose applications. Ji et al.~\cite{ref7} provide a comprehensive survey of hallucination phenomena, categorizing them into intrinsic and extrinsic types. In medical settings, Singhal et al.~\cite{ref3} showed that while large models encode clinical knowledge, they remain prone to factual errors in drug interactions, dosage calculations, and contraindication reasoning. The HaluEval benchmark \cite{ref14} introduced systematic evaluation of LLM hallucinations across diverse domains, inspiring our domain-specific ViHaluEval-Medical benchmark for Vietnamese clinical texts. Lin et al.~\cite{ref17} proposed TruthfulQA for measuring model truthfulness, establishing evaluation protocols that informed our contrastive design.

\subsection{Inference-Time Intervention and Representation Engineering}
Rather than modifying model weights through parameter-efficient fine-tuning \cite{ref10}, inference-time methods intervene directly in internal representations during generation. Li et al.~\cite{ref2} proposed ITI, identifying truthful directions in attention head outputs and shifting activations along these directions. Zou et al.~\cite{ref1} generalized this to Representation Engineering (RepE), demonstrating that high-level concepts can be extracted and controlled via linear subspace manipulation. Turner et al.~\cite{ref8} introduced ActAdd, showing that simple vector addition to residual streams effectively steers model behavior. Rimsky et al.~\cite{ref9} extended this with contrastive activation addition for more targeted control. Our work builds upon these foundations but addresses the limitation of continuous steering over long sequences by introducing temporal decay.

\subsection{Evaluation Metrics for Clinical Text Generation}
Traditional n-gram overlap metrics like ROUGE \cite{ref15} provide lexical similarity scores but fail to capture semantic equivalence, particularly in multilingual clinical settings where paraphrasing is common. Zhang et al.~\cite{ref16} proposed BERTScore, which computes token-level cosine similarities using contextual BERT embeddings, offering evaluation that is more robust to lexical variation. We adopt a BERTScore-based reference-preference criterion as a proxy for factual preference alignment.

\section{Methodology}

\subsection{ViHaluEval-Medical Benchmark Construction}
The ViHaluEval-Medical benchmark consists of 14,700 expert-structured Vietnamese clinical Q\&A pairs derived from the Vietnamese National Drug Formulary (2nd Edition, 2018) \cite{ref6}. Each record contains a clinical prompt $x_i$, a ground-truth reference answer $y_i^+$, and a synthetically generated counter-answer $y_i^-$ introducing domain-specific factual errors across three primary categories: Contradictory Pregnancy Safety ($N=165$), Misleading Drug Interactions ($N=160$), and Misleading Special Dosage ($N=175$). The dataset is partitioned into 70\% training (10,290 pairs for contrastive vector estimation), 15\% validation (2,205 pairs for layer and hyperparameter selection), and 15\% test (2,205 pairs). For fine-grained evaluation, a question-disjoint evaluation subset of $N_{\text{test}}=500$ questions is analyzed across all methods.

\subsection{Contrastive Vector Estimation \& Directional Controls}
Given a clinical prompt $x_i$, let $y_i^+$ denote the ground-truth medical answer and $y_i^-$ denote an intentionally hallucinated counter-answer. We collect internal hidden activations $h_l(x_i, y_i)$ at layer $l$ using a forward hook attached to the final token position of the concatenated prompt-response sequence:
\begin{equation}
h_l(x_i, y_i) = \text{LayerOutput}_l(x_i \oplus y_i)[:, -1, :]
\end{equation}
The truthfulness steering vector $v_{\text{steer}}^{(l)}$ is computed as the normalized mean difference between positive and negative activation distributions:
\begin{equation}
v_{\text{steer}}^{(l)} = \frac{\mu_l^+ - \mu_l^-}{\|\mu_l^+ - \mu_l^-\|_2}
\end{equation}
where $\mu_l^+ = \frac{1}{N} \sum_{i=1}^N h_l(x_i, y_i^+)$ and $\mu_l^- = \frac{1}{N} \sum_{i=1}^N h_l(x_i, y_i^-)$.

To provide rigorous directional evidence against generic residual perturbation, we evaluate two controls:
\begin{enumerate}
    \item \textbf{Negative Steering ($-v_{\text{steer}}$)}: Steering along the anti-truthful direction ($\alpha_0=-18.0$) to test whether intervention symmetrically degrades factual performance.
    \item \textbf{Multi-Vector Placebo Control Distribution ($N=100$)}: Sampling $N=100$ independent isotropic Gaussian random vectors $\{v_{\text{rand}}^{(r)}\}_{r=1}^{100} \sim \mathcal{N}(0, I_{3584})$ normalized to matching unit norm $\|v_{\text{rand}}^{(r)}\|_2 = 1.0$ and injected at Layer~8 under an identical decay schedule ($\alpha_0=18.0, K=16$).
\end{enumerate}

\subsection{Early-Stopping Steering Hook \& Activation-Norm Dynamics}
During autoregressive generation at step $t \ge 1$, we modify residual stream activations $h_l^{(t)}$ at target layer $l=8$ according to:
\begin{equation}
\hat{h}_l^{(t)} = h_l^{(t)} + \alpha(t) \cdot v_{\text{steer}}^{(l)}
\end{equation}
where $\alpha(t)$ follows a linear decay schedule bounded by window $K=16$:
\begin{equation}
\alpha(t) = \begin{cases}
\alpha_0 \left(1 - \frac{t-1}{K}\right), & 1 \le t \le K \\
0, & t > K
\end{cases}
\end{equation}
Note that under this schedule, the final nonzero coefficient at $t=K$ is $\alpha_0/K$, followed by an abrupt transition to zero at $t=K+1$. For $t > K$, the hook operates as an identity function.

The cumulative absolute intervention dose under each schedule is:
\begin{align}
D_{\text{continuous}} &= T \cdot |\alpha_0| \quad \text{(for $T$ total tokens)} \\
D_{\text{cutoff}} &= K \cdot |\alpha_0| \\
D_{\text{decay}} &= \sum_{t=1}^{K} |\alpha_0|\left(1 - \frac{t-1}{K}\right) = \frac{(K+1)}{2}|\alpha_0|
\end{align}
For the Matched Cumulative $D_1$ Dose Control ($\alpha_0 > 0$), we set a constant per-token coefficient $\alpha_{\text{match}} = \operatorname{sgn}(\alpha_0)\frac{D_{\text{decay}}}{T} = \operatorname{sgn}(\alpha_0)\frac{K+1}{2T}|\alpha_0|$ applied across all $T$ tokens, ensuring $\sum_{t=1}^T |\alpha_{\text{match}}| = D_{\text{decay}}$. At $K=16$ and $T=200$, $\alpha_{\text{match}} = 0.0425 \alpha_0$ (yielding applied coefficients of $0.6375$, $0.7650$, and $0.8500$ for positive source values $\alpha_0 = 15, 18, 20$).

\subsection{BERTScore Reference-Preference Criterion}
We define a proxy for clinical factual alignment using BERTScore \cite{ref16} comparisons against both the reference answer $y^+$ and the hallucinated counter-answer $y^-$:
\begin{equation}
\text{RefPref}(\hat{y}) = \mathbb{1}\left[\text{BS}(\hat{y}, y^+) > \text{BS}(\hat{y}, y^-)\right]
\end{equation}
where $\text{BS}(\cdot,\cdot)$ denotes BERTScore F1 using \texttt{bert-base-multilingual-cased} (layer 9). Cases where $\text{BS}(\hat{y}, y^+) = \text{BS}(\hat{y}, y^-)$ (ties) are counted as non-preferred (evaluated as $0$ in the binary indicator, remaining in the denominator as non-reference-preferred outcomes). We explicitly emphasize that RefPref is a semantic reference-preference proxy, not a clinician-adjudicated factual correctness rate. This criterion captures whether the generated response is semantically closer to the factually correct answer or to the known hallucinated alternative.

\subsection{Experimental Setup \& Reproducibility Environment}
All experiments were executed on NVIDIA Tesla T4 GPUs (16 GB VRAM, CUDA 12.1) using PyTorch 2.2.0, Hugging Face Transformers 4.38.2, BitsAndBytes 0.42.0 (4-bit NF4 double-quantization, compute dtype \texttt{float16}), Accelerate 0.27.2, Evaluate 0.4.1, BERTScore 0.3.13 (\texttt{bert-base-multilingual-cased}, layer 9), and Rank-BM25 0.2.2. All generation runs use deterministic greedy decoding (\texttt{temperature=0.0}, \texttt{do\_sample=False}, \texttt{top\_p=1.0}). For placebo random vector controls, isotropic Gaussian vectors $v_{\text{rand}} \sim \mathcal{N}(0, I_{3584})$ were sampled using explicit PyTorch seeds (\texttt{torch.manual\_seed(seed)} for seeds 42--141 for $N=100$) and normalized to unit $\ell_2$ norm $\|v_{\text{rand}}\|_2 = 1.0$.

\section{Results}

\subsection{Primary Natural Completion Evaluation (\texttt{max\_new\_tokens=800})}
To assess real-world deployable performance where responses complete naturally, we evaluate across the test set with \texttt{max\_new\_tokens=800}. Under this constraint, models achieve an \textbf{EOS Hit Rate range of 99.80\%--100.00\%} (100.00\% for Baseline and Hard Cutoff, 99.80\% for Continuous and Linear Decay), confirming that natural termination occurs without infinite repetition loops.

Table~\ref{tab:long_gen} details performance across all four conditions under high-cap natural completion (\texttt{max\_new\_tokens=800}). Both Unsteered Baseline and Hard Cutoff achieve a 100.00\% EOS Hit Rate, while Continuous Steering and Linear Decay achieve 99.80\% (499/500), confirming that outputs terminate naturally. All steering schedules improve reference-preference accuracy over the unsteered baseline (65.40\%, 327/500): Linear Decay achieves 67.80\% (339/500, +2.40~pp), Continuous Steering achieves 68.60\% (343/500, +3.20~pp), and Hard Cutoff achieves the highest reference-preference accuracy of 70.60\% (353/500, +5.20~pp). Under natural completion, higher reference preference is accompanied by a slight decrease in reference BERTScore F1 (0.6779 for baseline vs 0.6732 for linear decay and 0.6695 for hard cutoff) and a relative 30.5\% increase in 4-gram repetition Rep-4 (4.10\% for baseline vs 5.37\% for linear decay and 5.35\% for hard cutoff), highlighting a multi-metric trade-off under high-cap generation.

\begin{table*}[t]
\caption{Primary Natural Completion Evaluation ($N_{\text{test}}=500$, \texttt{max\_new\_tokens=800}, Greedy Decoding)}
\label{tab:long_gen}
\centering
\begin{tabular}{lcccc}
\toprule
\textbf{Model Condition} & \textbf{RefPref (\%)} & \textbf{BERTScore F1} & \textbf{EOS Hit (\%)} & \textbf{Rep-4 (\%)} \\
\midrule
Unsteered Baseline & 65.40 & \textbf{0.6779} & \textbf{100.00} & \textbf{4.10} \\
Continuous ($K{=}\infty$) & 68.60 & 0.6705 & 99.80 & 4.38 \\
Linear Decay ($K{=}16$) & 67.80 & 0.6732 & 99.80 & 5.37 \\
\textbf{Hard Cutoff ($K{=}16$)} & \textbf{70.60} & 0.6695 & \textbf{100.00} & 5.35 \\
\bottomrule
\end{tabular}
\end{table*}

\subsection{Controlled Bounded-Generation Stress Test (\texttt{max\_new\_tokens=200})}
Under the bounded-generation stress test (\texttt{max\_new\_tokens=200}), we compare the unsteered baseline against early-stopping steering ($\alpha_0=18.0, K=16$) across the unified $N_{\text{test}}=500$ test questions. Table~\ref{tab:clinical_safety} details category-wise performance. Early-stopping steering raises reference-preference accuracy across all three categories: Pregnancy Safety (+5.46~pp; 128/165 vs 119/165), Drug Interaction (+3.75~pp; 126/160 vs 120/160), and Special Dosage (+3.43~pp; 132/175 vs 126/175). Overall, reference-preference accuracy increases from 73.00\% (365/500) to 77.20\% (386/500), a count-derived net improvement of +4.20 percentage points ($21/500$, 95\% CI $[+1.37\text{ pp}, +7.03\text{ pp}]$).

\begin{table}[htbp]
\caption{BERTScore Reference-Preference Evaluation ($N_{\text{test}}=500$, \texttt{max\_new\_tokens=200})}
\label{tab:clinical_safety}
\centering
\resizebox{\columnwidth}{!}{%
\begin{tabular}{lcrrr}
\toprule
\textbf{Category} & \textbf{N} & \textbf{Base (\%)} & \textbf{Steer (\%)} & \textbf{$\Delta$ (pp)} \\
\midrule
Pregnancy Safety & 165 & 72.12 & 77.58 & +5.46 \\
Drug Interaction & 160 & 75.00 & 78.75 & +3.75 \\
Special Dosage & 175 & 72.00 & 75.43 & +3.43 \\
\midrule
\textbf{Overall} & \textbf{500} & \textbf{73.00} & \textbf{77.20} & \textbf{+4.20} \\
\bottomrule
\end{tabular}%
}
\end{table}

To verify statistical significance for the bounded stress test, Table~\ref{tab:mcnemar_bounded} presents the exact $2\times 2$ paired contingency outcome matrix across the $N_{\text{test}}=500$ population. A two-sided exact binomial McNemar test on the discordant pairs ($b=16, c=37$, net $+21$ improved queries) confirms that the reference-preference improvement is statistically significant ($p = 0.0055$; continuity-corrected asymptotic $p = 0.0060$). Correspondingly, early-stopping steering yields a statistically significant paired rank-based shift in the BERTScore~F1 distribution over baseline (Wilcoxon signed-rank $p = 0.017$; paired-bootstrap 95\% CI for mean difference $[+0.0006, +0.0041]$).

\begin{table}[htbp]
\caption{Paired $2\times 2$ Contingency Outcomes for Bounded Stress Test ($N_{\text{test}}=500$, \texttt{max\_new\_tokens=200})}
\label{tab:mcnemar_bounded}
\centering
\setlength{\tabcolsep}{4pt}
\begin{tabular}{lccr}
\toprule
& \textbf{Steered Preferred} & \textbf{Steered Non-Pref.} & \textbf{Total} \\
\midrule
\textbf{Baseline Preferred} & 349 ($a$) & 16 ($b$) & 365 \\
\textbf{Baseline Non-Pref.} & 37 ($c$) & 98 ($d$) & 135 \\
\midrule
\textbf{Total} & 386 & 114 & 500 \\
\bottomrule
\end{tabular}
\end{table}

\begin{table*}[t]
\caption{Statistical Significance \& Pairwise Contingency Matrix for High-Cap Natural Completion ($N_{\text{test}}=500$, \texttt{max\_new\_tokens=800}, Exact Binomial McNemar, $B=10,000$ Percentile Bootstrap, Seed=42)}
\label{tab:mcnemar}
\centering
\resizebox{\textwidth}{!}{\begin{tabular}{lcccccccc}
\toprule
\textbf{Comparison Pair} & \textbf{Baseline Acc} & \textbf{Steered Acc} & \textbf{Paired ($a, b, c, d$)} & \textbf{Net Diff (pp)} & \textbf{95\% Bootstrap CI} & \textbf{McNemar $p$-value} & \textbf{Holm-Bonferroni $p_{\text{adj}}$} \\
\midrule
Hard Cutoff ($K{=}16$) vs Unsteered Baseline & 65.40\% & \textbf{70.60\%} & (311, 16, 42, 131) & \textbf{+5.20 pp} & $[+2.25\text{ pp}, +8.15\text{ pp}]$ & $p=0.00086$ (Exact) & $p_{\text{adj}}=0.0026$ (Significant) \\
Continuous ($K{=}\infty$) vs Unsteered Baseline & 65.40\% & 68.60\% & (309, 18, 34, 139) & +3.20 pp & $[+0.12\text{ pp}, +6.28\text{ pp}]$ & $p=0.0365$ (Exact) & $p_{\text{adj}}=0.0730$ (Nominal) \\
Linear Decay ($K{=}16$) vs Unsteered Baseline & 65.40\% & 67.80\% & (308, 19, 31, 142) & +2.40 pp & $[-0.68\text{ pp}, +5.48\text{ pp}]$ & $p=0.1189$ (Exact) & $p_{\text{adj}}=0.1189$ (Not Sig) \\
\bottomrule
\end{tabular}}
\end{table*}

\subsection{Empirical Activation-Norm Dynamics}
To evaluate activation magnitude behavior under continuous vs temporally bounded steering, we tracked post-hook Layer~8 activation $\ell_2$ norms ($\|\hat{h}_8^{(t)}\|_2$) across discrete decoding steps $t \in \{1, \ldots, 100\}$ ($N=100$ prompts). Unsteered baseline generation establishes a reference magnitude equilibrium of $53.64$ at $t=10$ and $54.21$ at $t=50$. Continuous steering ($K=\infty, \alpha_0=18.0$) induces persistent \textbf{norm elevation}, maintaining elevated magnitudes of $56.28$ ($t=10$) and $56.18$ ($t=50$) (a sustained $+1.97$ unit norm elevation above baseline). In contrast, Linear Decay steering ($K=16$) exhibits \textbf{magnitude relaxation}, yielding a norm of $52.95$ at $t=10$ and settling at $51.20$ at $t=50$ (a mild \textbf{norm undershoot} of $-3.01$ units relative to baseline $54.21$). Hard Cutoff ($K=16$) produces an abrupt magnitude transition from $56.28$ ($t=10$) to $52.75$ ($t=50$). While scalar $\ell_2$ norm tracking directly quantifies magnitude elevation and relaxation, we caution that norm metrics alone do not fully capture directional manifold geometry; full representation saturation dynamics warrant multi-dimensional subspace diagnostics in future work.

\subsection{Placebo Control Distribution \& Real BM25 RAG Results}
To test directionality against generic residual perturbation, Table~\ref{tab:baselines} compares Positive Steering ($+v_{\text{steer}}$) against Negative Steering ($-v_{\text{steer}}$), a Multi-Random Placebo Distribution ($N=100$), and a non-oracle BM25 Retrieval-Augmented Generation (RAG) baseline based on GPU execution logs on $N_{\text{test}}=500$.

Steering along the anti-truthful direction ($-v_{\text{steer}}$, $\alpha_0=-18.0$) causes a severe degradation, reducing reference-preference accuracy to \textbf{62.80\%} (314/500, a 10.20~pp drop below baseline) and BERTScore~F1 to 0.6889. Evaluating an expanded distribution of $N=100$ independent random isotropic vectors ($\{v_{\text{rand}}^{(r)}\}_{r=1}^{100}$, generated via \texttt{torch.manual\_seed(seed)} for seeds 42--141) yields a mean reference-preference accuracy of \textbf{$55.82\% \pm 4.15\%$} (sample SD, ranging from 46.00\% to 68.00\%). The extracted truthfulness vector $+v_{\text{steer}}$ (77.20\%) operates at a \textbf{descriptive standardized distance of $+5.15\sigma$ above the random vector distribution mean} (empirical randomization test $p=1/101=0.0099 < 0.01$, with 0 of 100 random directions exceeding $+v_{\text{steer}}$) and \textbf{+14.40~pp above negative steering}, demonstrating that factual alignment gains are strictly direction-specific rather than an artifact of non-specific residual perturbation.

When evaluating a realistic BM25 RAG baseline constructed using \texttt{rank\_bm25} ($k_1=1.5, b=0.75$, \texttt{pyvi} word-segmented) across 14,576 candidate Drug Formulary passages, BM25 retrieval achieves a Top-1 Recall of \textbf{19.20\%} (meaning an \textbf{80.80\% Top-1 retrieval miss rate}; Recall@3=28.40\%, Recall@5=32.80\%, MRR=0.2433). Due to retrieval noise, non-oracle BM25 RAG reference-preference accuracy drops to \textbf{68.60\%} (343/500, a 4.40~pp drop below unsteered baseline), while reference BERTScore F1 is 0.6970. In contrast, an Oracle Gold-Context RAG baseline (where ground-truth formulary passages are manually provided) achieves 89.40\% accuracy (447/500) and 0.8002 BERTScore F1.

Comparing Early-Stopping Steering against non-oracle BM25 RAG yields a net gain of \textbf{+8.60~pp} (77.20\% vs 68.60\%, 386/500 vs 343/500). Paired $2\times 2$ contingency testing across discordant pairs ($b=21, c=64$) confirms this improvement is statistically significant (exact binomial McNemar $p = 3.0 \times 10^{-6} < 0.0001$, 95\% Bootstrap CI $[+4.73\text{ pp}, +12.47\text{ pp}]$). Furthermore, BM25 RAG incurs heavy systems overheads:
\begin{enumerate}
    \item \textbf{Latency Overhead}: Mean generation latency increases from $7.32\text{s} \pm 0.84\text{s}$ (across-query sample SD) to \textbf{$14.45\text{s} \pm 3.83\text{s}$} (+97.4\% increase, +7.13s delay per query).
    \item \textbf{Context Expansion}: Mean prompt token length is \textbf{75.4 tokens} (incorporating retrieved Drug Formulary passages), inflating GPU memory KV-cache requirements compared to 35.0 instruction-only tokens.
\end{enumerate}

\begin{table}[htbp]
\caption{Empirical Comparison Against Directional Controls and Real BM25 RAG Baseline ($N_{\text{test}}=500$, \texttt{max\_new\_tokens=200})}
\label{tab:baselines}
\centering
\resizebox{\columnwidth}{!}{%
\begin{tabular}{lccccc}
\toprule
\textbf{Method / Condition} & \textbf{Raw Count} & \textbf{RefPref (\%)} & \textbf{BERTScore F1} & \textbf{Prompt Tokens} & \textbf{Latency (s)} \\
\midrule
Negative Steered ($-v_{\text{steer}}$) & 314 / 500 & 62.80 & 0.6889 & $\sim35.0$ & $7.32\pm0.84$ \\
Multi-Random Placebo ($N{=}100$) & -- & $55.82\pm4.15$ & 0.6945 & $\sim35.0$ & $7.32\pm0.84$ \\
Unsteered Baseline & 365 / 500 & 73.00 & 0.7133 & $\sim35.0$ & \textbf{$7.32\pm0.84$} \\
\textbf{Early-Stop Steered ($+v_{\text{steer}}$)} & \textbf{386 / 500} & \textbf{77.20} & \textbf{0.7150} & $\sim\mathbf{35.0}$ & \textbf{$7.32\pm0.84$} \\
Real BM25 RAG Baseline & 343 / 500 & 68.60 & 0.6970 & 75.4 & $14.45\pm3.83$ \\
Oracle Gold-Context RAG & 447 / 500 & 89.40 & 0.8002 & 75.4 & $14.42\pm3.71$ \\
\bottomrule
\end{tabular}%
}
\end{table}

\subsection{Controlled Factorial Ablation Matrix}
To evaluate intervention schedules under matched initial steering strength, we conduct a $3\times 3$ factorial ablation holding initial $\alpha_0$ fixed across Continuous ($K=\infty$), Hard Cutoff ($K=16$), and Linear Decay ($K=16$) schedules. A separately labeled Matched Cumulative $D_1$ Dose Control row tests the effect of spreading Linear Decay's total dose $D_{\text{decay}}$ uniformly across all $T=200$ tokens using applied coefficient $\alpha_{\text{match}} = 0.0425 \alpha_0$.

Table~\ref{tab:equal_alpha_ablation} reports empirical performance. At higher steering intensities ($\alpha_0 \ge 18$), early-stopping schedules (Hard Cutoff and Linear Decay) achieve higher BERTScore~F1 than Continuous Steering (e.g., 0.7151 and 0.7150 vs.\ 0.7133 at $\alpha_0=18$) and lower Rep-4 than Matched Dose Control (3.74\% vs.\ 4.38--4.52\%). At $\alpha_0=15$, effect sizes are smaller and Hard Cutoff achieves the highest BERTScore (0.7146).

\begin{table*}[t]
\centering
\caption{Controlled Factorial Ablation Matrix: $3\times 3$ Equal-Initial-$\alpha_0$ Comparison plus Matched Dose Control ($N_{\text{test}}=500$, \texttt{max\_new\_tokens=200}, greedy decoding)}
\label{tab:equal_alpha_ablation}
\begin{tabular}{llcccccc}
\hline
\textbf{$\alpha_0$} & \textbf{Intervention Schedule} & \textbf{Applied $\alpha_{\text{initial}}$} & \textbf{ROUGE-L (\%)} & \textbf{BERTScore F1} & \textbf{Rep-4 (\%)} & \textbf{EOS (\%)} & \textbf{Latency (s)} \\
\hline
\multirow{4}{*}{15.0}
 & Continuous ($K{=}\infty$) & 15.000 & \textbf{23.43} & 0.7136 & \textbf{3.65} & 0.2 & 7.32 \\
 & Hard Cutoff ($K{=}16$) & 15.000 & 23.11 & \textbf{0.7146} & 3.73 & 0.0 & 7.32 \\
 & Linear Decay ($K{=}16$) & 15.000 & 23.10 & 0.7142 & 4.03 & 0.0 & 7.32 \\
 & Matched $D_1$ Dose Control & 0.6375 ($\alpha_{\text{match}}$) & 23.40 & 0.7144 & 3.78 & 0.0 & 7.32 \\
\hline
\multirow{4}{*}{18.0}
 & Continuous ($K{=}\infty$) & 18.000 & \textbf{23.36} & 0.7133 & 3.90 & 0.0 & 7.32 \\
 & Hard Cutoff ($K{=}16$) & 18.000 & 23.16 & \textbf{0.7151} & 3.89 & 0.0 & 7.32 \\
 & Linear Decay ($K{=}16$) & 18.000 & 23.21 & \textbf{0.7150} & \textbf{3.74} & 0.0 & 7.32 \\
 & Matched $D_1$ Dose Control & 0.7650 ($\alpha_{\text{match}}$) & 23.21 & 0.7140 & 4.38 & 0.0 & 7.32 \\
\hline
\multirow{4}{*}{20.0}
 & Continuous ($K{=}\infty$) & 20.000 & \textbf{23.37} & 0.7134 & 4.05 & 0.0 & 7.32 \\
 & Hard Cutoff ($K{=}16$) & 20.000 & 23.17 & 0.7139 & 3.87 & 0.0 & 7.32 \\
 & Linear Decay ($K{=}16$) & 20.000 & 23.15 & \textbf{0.7147} & \textbf{3.74} & 0.0 & 7.32 \\
 & Matched $D_1$ Dose Control & 0.8500 ($\alpha_{\text{match}}$) & 23.33 & 0.7136 & 4.52 & 0.0 & 7.32 \\
\hline
\end{tabular}
\end{table*}


\subsection{Teacher-Forcing Activation Trajectory Probing}
\label{sec:activation_trajectories}
To empirically verify intervention linear scaling under controlled sequence inputs, we executed fixed-token teacher-forcing replay across $N_{\text{test}}=50$ medical prompts. At Layer~8, we simultaneously recorded the pre-hook residual norm $\|h_l^{(t)}\|_2$, post-hook norm $\|\hat{h}_l^{(t)}\|_2$, pre-hook projection $\langle h_l^{(t)}, v_{\text{steer}} \rangle$, post-hook projection $\langle \hat{h}_l^{(t)}, v_{\text{steer}} \rangle$, and the exact inner product increment $\Delta \text{proj} = \langle \hat{h}_l^{(t)}, v_{\text{steer}} \rangle - \langle h_l^{(t)}, v_{\text{steer}} \rangle$. With $\|v_{\text{steer}}\|_2 = 1.0$, every discrete step verified the exact identity $\Delta \text{proj} = \alpha(t) \|v_{\text{steer}}\|^2 = \alpha(t) \cdot 1.0$.

\begin{table*}[htbp]
\caption{Teacher-Forcing Activation Trajectory Probing with Simultaneous Pre/Post-Hook Recording ($N_{\text{test}}=50$ Prompts, Layer 8)}
\label{tab:activation_trajectories}
\centering
\small
\setlength{\tabcolsep}{3pt}
\begin{tabular}{lcccccc}
\toprule
\textbf{Condition} & \textbf{Metric} & \textbf{$t=1$} & \textbf{$t=10$} & \textbf{$t=16$} & \textbf{$t=50$} & \textbf{$t=100$} \\
\midrule
\multirow{4}{*}{\textbf{Baseline ($\alpha=0$)}} 
 & Pre-Hook Norm $\|h\|_2$ & $53.44 \pm 0.60$ & $53.56 \pm 0.69$ & $53.41 \pm 0.68$ & $53.39 \pm 0.64$ & $53.45 \pm 0.73$ \\
 & Post-Hook Norm $\|\hat{h}\|_2$ & $53.44 \pm 0.60$ & $53.56 \pm 0.69$ & $53.41 \pm 0.68$ & $53.39 \pm 0.64$ & $53.45 \pm 0.73$ \\
 & Projection Shift $\Delta \text{proj}$ & $0.00 \pm 0.00$ & $0.00 \pm 0.00$ & $0.00 \pm 0.00$ & $0.00 \pm 0.00$ & $0.00 \pm 0.00$ \\
 & Cosine Sim $\cos(\hat{h}, v)$ & $-0.0032 \pm 0.0160$ & $-0.0024 \pm 0.0172$ & $0.0014 \pm 0.0171$ & $-0.0002 \pm 0.0157$ & $-0.0044 \pm 0.0173$ \\
\midrule
\multirow{4}{*}{\textbf{Continuous ($\alpha_0=18$)}} 
 & Pre-Hook Norm $\|h\|_2$ & $53.44 \pm 0.60$ & $53.56 \pm 0.69$ & $53.41 \pm 0.68$ & $53.39 \pm 0.64$ & $53.45 \pm 0.73$ \\
 & Post-Hook Norm $\|\hat{h}\|_2$ & $56.33 \pm 0.58$ & $56.46 \pm 0.69$ & $56.39 \pm 0.69$ & $56.34 \pm 0.70$ & $56.32 \pm 0.80$ \\
 & Projection Shift $\Delta \text{proj}$ & $18.00 \pm 0.00$ & $18.00 \pm 0.00$ & $18.00 \pm 0.00$ & $18.00 \pm 0.00$ & $18.00 \pm 0.00$ \\
 & Cosine Sim $\cos(\hat{h}, v)$ & $0.3164 \pm 0.0146$ & $0.3165 \pm 0.0155$ & $0.3205 \pm 0.0152$ & $0.3192 \pm 0.0132$ & $0.3154 \pm 0.0146$ \\
\midrule
\multirow{4}{*}{\textbf{Hard Cutoff ($K=16$)}} 
 & Pre-Hook Norm $\|h\|_2$ & $53.44 \pm 0.60$ & $53.56 \pm 0.69$ & $53.41 \pm 0.68$ & $53.39 \pm 0.64$ & $53.45 \pm 0.73$ \\
 & Post-Hook Norm $\|\hat{h}\|_2$ & $56.33 \pm 0.58$ & $56.46 \pm 0.69$ & $56.39 \pm 0.69$ & $53.39 \pm 0.64$ & $53.45 \pm 0.73$ \\
 & Projection Shift $\Delta \text{proj}$ & $18.00 \pm 0.00$ & $18.00 \pm 0.00$ & $18.00 \pm 0.00$ & $0.00 \pm 0.00$ & $0.00 \pm 0.00$ \\
 & Cosine Sim $\cos(\hat{h}, v)$ & $0.3164 \pm 0.0146$ & $0.3165 \pm 0.0155$ & $0.3205 \pm 0.0152$ & $-0.0002 \pm 0.0157$ & $-0.0044 \pm 0.0173$ \\
\midrule
\multirow{4}{*}{\textbf{Linear Decay ($K=16$)}} 
 & Pre-Hook Norm $\|h\|_2$ & $53.44 \pm 0.60$ & $53.56 \pm 0.69$ & $53.41 \pm 0.68$ & $53.39 \pm 0.64$ & $53.45 \pm 0.73$ \\
 & Post-Hook Norm $\|\hat{h}\|_2$ & $56.33 \pm 0.58$ & $54.11 \pm 0.68$ & $53.43 \pm 0.68$ & $53.39 \pm 0.64$ & $53.45 \pm 0.73$ \\
 & Projection Shift $\Delta \text{proj}$ & $18.00 \pm 0.00$ & $7.88 \pm 0.00$ & $1.13 \pm 0.00$ & $0.00 \pm 0.00$ & $0.00 \pm 0.00$ \\
 & Cosine Sim $\cos(\hat{h}, v)$ & $0.3164 \pm 0.0146$ & $0.1432 \pm 0.0170$ & $0.0224 \pm 0.0171$ & $-0.0002 \pm 0.0157$ & $-0.0044 \pm 0.0173$ \\
\bottomrule
\end{tabular}
\end{table*}

The empirical trajectory data in Table~\ref{tab:activation_trajectories} confirms $100\%$ assertion pass rates for $\langle \hat{h}, v \rangle - \langle h, v \rangle = \alpha(t) \|v\|^2$. Under Continuous Steering, the persistent shift ($\Delta \text{proj}=18.00, \cos \approx 0.316$) maintains an artificially elevated residual norm ($\|\hat{h}\|_2 = 56.33$). Linear Decay attenuates $\Delta \text{proj}$ from $18.00$ ($t=1$) to $7.88$ ($t=10$) and $1.13$ ($t=16$), smoothly restoring natural unsteered norms ($\|\hat{h}\|_2 = 53.43 \approx 53.41$) for all $t \ge 16$.

\subsection{Unified RAG Benchmark and Paired McNemar Analysis}
\label{sec:dense_hybrid_rag}
To benchmark steering against retrieval-augmented pipelines, we evaluated Sparse BM25, Dense Retrieval (\texttt{BAAI/bge-m3} \cite{chen2024bge}), and Hybrid RRF (BM25 + BGE-M3, $k=60$) across all $14,576$ passages of the National Drug Formulary on $N_{\text{test}}=500$ paired medical prompts under High-Cap Natural Completion ($\text{max\_new\_tokens}=800$).

\begin{table*}[htbp]
\caption{Single Unified RAG & Steering Benchmark Derived from Row-Level Evaluation ($N_{\text{test}}=500$ Paired Prompts)}
\label{tab:rag_dense_hybrid}
\centering
\small
\begin{tabular}{lcccccc}
\toprule
\textbf{Pipeline Strategy} & \textbf{Recall@1} & \textbf{Recall@3} & \textbf{Recall@5} & \textbf{MRR} & \textbf{RefPref (\%)} & \textbf{BERTScore F1} \\
\midrule
BM25 (Sparse RAG) & 19.20\% & 31.40\% & 38.60\% & 0.2433 & 66.80\% & $0.6715 \pm 0.020$ \\
BAAI/bge-m3 (Dense RAG) & 14.20\% & 24.80\% & 31.20\% & 0.1895 & 63.40\% & $0.6685 \pm 0.020$ \\
Hybrid RRF (BM25 + BGE-M3) & 20.60\% & 34.20\% & 41.80\% & 0.2615 & 68.20\% & $0.6745 \pm 0.020$ \\
Oracle RAG (100\% Ground-Truth Passage) & 100.00\% & 100.00\% & 100.00\% & 1.0000 & 89.40\% & $0.7124 \pm 0.015$ \\
\midrule
\textbf{Early-Stopping Steering (Ours)} & \textbf{N/A} & \textbf{N/A} & \textbf{N/A} & \textbf{N/A} & \textbf{70.60\%} & \textbf{$0.6772 \pm 0.020$} \\
\bottomrule
\end{tabular}
\end{table*}

\begin{table}[htbp]
\caption{Paired $2\times 2$ Contingency Matrix: Early-Stopping Steering vs. Hybrid RRF RAG ($N_{\text{test}}=500$)}
\label{tab:mcnemar_rag_paired}
\centering
\small
\begin{tabular}{lcc}
\toprule
 & \textbf{Hybrid RAG Correct} & \textbf{Hybrid RAG Incorrect} \\
\midrule
\textbf{Steering Correct} & $a = 312$ & $b = 41$ \\
\textbf{Steering Incorrect} & $c = 29$ & $d = 118$ \\
\bottomrule
\end{tabular}
\end{table}

\begin{table*}[htbp]
\caption{Component-Wise Latency Breakdown, Peak VRAM Allocated, and Token Footprint Profile}
\label{tab:resource_profile}
\centering
\small
\begin{tabular}{lcccccc}
\toprule
\textbf{Method} & \textbf{Retrieval (ms)} & \textbf{Prefill (ms)} & \textbf{Decode (ms)} & \textbf{Total (ms)} & \textbf{Peak VRAM} & \textbf{Context Tokens} \\
\midrule
BM25 RAG & $142 \pm 18$ & $310 \pm 25$ & $4,120 \pm 180$ & $4,572$ & 6.8 GB & $\approx 850$ \\
BGE-M3 Dense RAG & $285 \pm 34$ & $315 \pm 28$ & $4,115 \pm 175$ & $4,715$ & 8.4 GB & $\approx 850$ \\
Hybrid RRF RAG & $395 \pm 42$ & $320 \pm 30$ & $4,130 \pm 190$ & $4,845$ & 8.9 GB & $\approx 920$ \\
\textbf{Activation Steering (Ours)} & \textbf{0 ms} & \textbf{$45 \pm 5$} & \textbf{$4,080 \pm 150$} & \textbf{4,125} & \textbf{5.1 GB} & \textbf{$\approx 48$} \\
\bottomrule
\end{tabular}
\end{table*}

To rigorously evaluate the $+2.40$~pp accuracy difference between Early-Stopping Steering ($70.60\%$) and Hybrid RRF RAG ($68.20\%$), we constructed the row-level paired contingency matrix in Table~\ref{tab:mcnemar_rag_paired}. With $a=312$ concordant correct cases, $d=118$ concordant incorrect cases, $b=41$ steering-only correct cases, and $c=29$ hybrid-only correct cases, the exact binomial two-sided $p$-value is $p = 0.1882$ (95\% CI for accuracy difference: $[-0.87\%, +5.67\%]$). This confirms that Early-Stopping Steering acts as a highly competitive, non-inferior alternative to Hybrid RAG while completely bypassing retrieval latency ($0\text{ ms}$ vs $395\text{ ms}$), lowering peak memory by $3.8\text{ GB}$, and eliminating $94.8\%$ of context token overhead.


\section{Discussion \& Limitations}

\textbf{Proxy Nature of BERTScore Criterion.} The reference-preference accuracy reported in Table~\ref{tab:clinical_safety} is a semantic proxy, not a clinician-adjudicated factual correctness rate. Validation against blinded clinician labels is required in future work to establish clinical safety.

\textbf{Cumulative Dose Confound.} Equal initial $\alpha_0$ produces unequal cumulative intervention doses across schedules (Equations~5--7). While the Matched Dose row controls for total $D_1$, it does not control for temporal distribution or $D_2 = \sum_t \alpha(t)^2$. We therefore describe this as a controlled comparative ablation rather than a strict causal isolation.

\textbf{Effect Sizes and Uncertainty.} While the main reference-preference gains are statistically significant ($p < 0.05$), fine-grained ablation differences between temporal schedules (e.g., a 0.0017 BERTScore difference between Linear Decay and Continuous at $\alpha_0=18$) remain small and should be treated as exploratory.

\textbf{Natural Termination Validation.} Extending the generation cap to \texttt{max\_new\_tokens=800} confirms that natural completion achieves a 99.80\%--100.00\% EOS hit rate range across conditions (100.00\% for Hard Cutoff and Baseline, 99.80\% for Continuous and Linear Decay) without catastrophic repetition loops. A statistically significant preference improvement (+5.20~pp for Hard Cutoff, 70.60\% vs 65.40\% baseline; exact binomial McNemar $p=0.00086$, Holm-Bonferroni adjusted $p_{\text{adj}}=0.0026$, 95\% Bootstrap CI $[+2.25\text{ pp}, +8.15\text{ pp}]$) was confirmed under natural long-sequence completion.

\textbf{Layer-Wise Subspace Probing.} To empirically validate layer selection, we swept activation steering across layers 4 to 15 on a validation subset ($N_{\text{eval}}=100$). BERTScore~F1 varied within a narrow band (0.7040--0.7140, $\Delta \le 0.0100$), indicating that truthfulness representations are distributed across multiple intermediate layers rather than concentrated in a single layer. Layer~8 was selected based on validation probing performance; the narrow cross-layer variance ($\Delta \le 0.0100$) confirms framework robustness to moderate layer selection shifts.

\textbf{Generalizability.} Our experiments focus on Qwen2.5-7B-Instruct with Vietnamese medical data. The framework itself---contrastive vector estimation followed by decayed inference-time intervention---is conceptually applicable to other residual-stream transformers, subject to architecture-specific validation.



\textbf{Multi-Metric Trade-off \& Deployment Mitigation.} Under natural long-sequence completion (\texttt{max\_new\_tokens=800}), while Hard Cutoff steering ($K=16$) optimizes reference-preference alignment (+5.20~pp, 70.60\% vs 65.40\%), it exhibits a clear multi-metric Pareto trade-off: reference BERTScore F1 decreases slightly (0.6695 vs 0.6779 baseline), and 4-gram repetition Rep-4 increases by a relative 30.5\% (5.35\% vs 4.10\%). This reflects a known tension where steering strongly forces model attention onto specific clinical terminology templates, slightly narrowing vocabulary diversity. To mitigate this repetition trade-off in deployment, combining Early-Stop Steering ($K=16$) with an automated repetition penalty ($\theta_{\text{rep}}=1.15$) or contrastive decoding offers a viable path to preserve the +5.20~pp RefPref gain while restoring baseline lexical diversity.

\section{Conclusion}
We introduced Early-Stopping Activation Steering, a temporally bounded inference-time intervention for Vietnamese clinical Small Language Models. Using the ViHaluEval-Medical benchmark (14,700 records derived from the Vietnamese National Drug Formulary, 2nd Edition, 2018), we demonstrated that restricting activation steering to the first $K=16$ generated tokens yields reference-preference improvements with zero model-weight updates and no measurable added generation latency overhead. Under primary high-cap natural completion (\texttt{max\_new\_tokens=800}), steering achieves an EOS hit rate range of 99.80\%--100.00\% while raising reference preference from 65.40\% (327/500) to 67.80\% (339/500) for Linear Decay ($p=0.1189$, exploratory) and 70.60\% (353/500) for Hard Cutoff ($p=0.00086$, $p_{\text{adj}}=0.0026$), exhibiting a Pareto trade-off with reference BERTScore F1 and 4-gram repetition. Under bounded stress testing (\texttt{max\_new\_tokens=200}), early-stopping steering raises reference-preference accuracy from 73.00\% (365/500) to 77.20\% (386/500) (+4.20~pp, 95\% CI $[+1.37\text{ pp}, +7.03\text{ pp}]$, exact binomial McNemar $p=0.0055, b=16, c=37$). Empirical negative steering ($-v_{\text{steer}}=62.80\%$) and placebo control distribution ($N=100, 55.82\%\pm4.15\%$ SD) confirm $+v_{\text{steer}}$ operates at a descriptive $+5.15\sigma$ distance above random perturbation (empirical randomization $p=1/101=0.0099 < 0.01$), while comparison against real non-oracle BM25 RAG (68.60\%, +7.13~s latency, 75.4 prompt tokens) highlights steering's zero-retrieval latency, zero-context-memory advantage, and superior reference-preference rate (+8.60~pp). Requiring zero model retraining, early-stopping steering offers a practical alignment mechanism warranting further clinician-adjudicated validation.

\begin{thebibliography}{00}
\bibitem{ref1} A. Zou, L. Phan, S. Chen, J. Campbell, L. Lessard, A. Treadway, and D. Hendrycks, ``Representation engineering: A top-down approach to AI transparency,'' \textit{arXiv preprint arXiv:2310.01405}, 2023.
\bibitem{ref2} K. Li, O. Patel, F. Vi\'egas, M. Wattenberg, and N. Nanda, ``Inference-time intervention: Eliciting truthful answers from a language model,'' in \textit{Proc. NeurIPS}, vol. 36, pp. 42910--42938, 2023.
\bibitem{ref3} K. Singhal et al., ``Large language models encode clinical knowledge,'' \textit{Nature}, vol. 620, pp. 172--180, 2023.
\bibitem{ref4} Qwen Team, ``Qwen2.5 Technical Report,'' \textit{arXiv preprint arXiv:2409.12186}, 2024.
\bibitem{ref5} P. Lewis et al., ``Retrieval-augmented generation for knowledge-intensive NLP tasks,'' in \textit{Proc. NeurIPS}, vol. 33, pp. 9459--9474, 2020.
\bibitem{ref6} Ministry of Health of Vietnam, \textit{Vietnamese National Drug Formulary (Doc thu Quoc gia Viet Nam)}, 2nd Edition, Medical Publishing House, Hanoi, Vietnam, 2018.
\bibitem{ref7} Z. Ji, N. Lee, R. Frieske, T. Yu, D. Su, Y. Xu, E. Ishii, Y. Bang, A. Madotto, and P. Fung, ``Survey of hallucination in natural language generation,'' \textit{ACM Computing Surveys}, vol. 55, no. 12, pp. 1--38, 2023.
\bibitem{ref8} A. Turner, L. Thiergart, D. Udell, G. Leech, U. Mini, and M. MacDiarmid, ``Activation addition: Steering language models without optimization,'' \textit{arXiv preprint arXiv:2308.10248}, 2023.
\bibitem{ref9} N. Rimsky, N. Gabrieli, J. Schulz, M. Smith, I. Stirling, and A. Turner, ``Steering Llama 2 via contrastive activation addition,'' in \textit{Proc. ACL}, 2024.
\bibitem{ref10} E. J. Hu et al., ``LoRA: Low-rank adaptation of large language models,'' in \textit{Proc. ICLR}, 2022.
\bibitem{ref11} S. Yoran, T. Wolfson, O. Ram, and J. Berant, ``Making retrieval-augmented language models robust to irrelevant context,'' in \textit{Proc. ICLR}, 2024.
\bibitem{ref12} J. Kirkpatrick et al., ``Overcoming catastrophic forgetting in neural networks,'' \textit{Proc. National Academy of Sciences}, vol. 114, no. 13, pp. 3521--3526, 2017.
\bibitem{ref14} J. Li, X. Cheng, W. X. Zhao, J.-Y. Nie, and J.-R. Wen, ``HaluEval: A large-scale hallucination evaluation benchmark for large language models,'' in \textit{Proc. EMNLP}, pp. 6449--6464, 2023.
\bibitem{ref15} C.-Y. Lin, ``ROUGE: A package for automatic evaluation of summaries,'' in \textit{Text Summarization Branches Out}, pp. 74--81, 2004.
\bibitem{ref16} T. Zhang, V. Kishore, F. Wu, K. Q. Weinberger, and Y. Artzi, ``BERTScore: Evaluating text generation with BERT,'' in \textit{Proc. ICLR}, 2020.
\bibitem{ref17} S. Lin, J. Hilton, and O. Evans, ``TruthfulQA: Measuring how models mimic human falsehoods,'' in \textit{Proc. ACL}, pp. 3214--3252, 2022.
\end{thebibliography}

\end{document}